\chapter{Main theorem: sparse spanners for temporal cliques}
\label{chap:main-theorem}

This is the target result of the project, now fully formalized. It was first
proved by Casteigts, Peters \& Schoeters, \emph{Temporal Cliques Admit Sparse
Spanners}, ICALP 2019 / JCSS. The formal proof follows the dismountability-only
proof of Daniele Carnevale, Arnaud Casteigts \& Timothée Corsini,
\emph{Dismountability in Temporal Cliques Revisited}, SAND 2025 (LIPIcs 330,
6:1--6:18, doi:10.4230/LIPIcs.SAND.2025.6, arXiv:2502.01321), \S\S~2--4, which
concludes with the one-sided dismountability of Angrick et al.\ (ESA 2024).

\begin{theorem}[$O(n \log n)$ spanners of temporal cliques --- SAND \S~4.1]
  \label{thm:main-spanner-bound}
  \uses{thm:clique-spanner-ncard-le,thm:spanner-bound-isBigO}
  \lean{TemporalGraph.Subgraph.exists_isBigO_isSpannerOf}
  \leanok
  There is a function $f(n) = O(n \log n)$, chosen once for all temporal
  cliques, such that every finite, simple, proper temporal clique $H$ on $n$
  vertices (with $H$ and its ambient graph symmetric, the ambient graph
  loopless) admits a symmetric spanner $K$ of $H$ (Definition~\ref{def:spanner})
  with at most $f(n)$ edges. One can take
  $f(n) = 2n\lceil\log_2 n\rceil + 4n$
  (Theorem~\ref{thm:clique-spanner-ncard-le}).
\end{theorem}

\begin{proof}
  \leanok
  Take $f(n) = 2n\lceil\log_2 n\rceil + 4n$, which bounds the spanners of
  every such clique by Theorem~\ref{thm:clique-spanner-ncard-le}. Since
  $\lceil\log_2 n\rceil \le \log_2 n + 1$ and $1 \le \log_2 n$ for $n \ge 2$,
  $f(n) \le 8 n \log_2 n$ eventually, so $f(n) = O(n \log n)$.
\end{proof}

\section*{Proof strategy}

Theorem~\ref{thm:clique-spanner-ncard-le} is proved by strong induction on the
number $n$ of vertices.
\begin{enumerate}
\item
  If $H$ is $\{1,2\}$-hop dismountable (Definition~\ref{def:two-hop-dismountable}),
  some vertex $x$ can be dismounted: a spanner of $H - x$, given by induction,
  extends to a spanner of $H$ with at most $4$ more edges
  (Theorem~\ref{thm:spanner-of-two-hop-dismountable}).
\item
  Otherwise $V^0 = \emptyset$ (Lemma~\ref{lem:vzero-empty}) and
  $|V^-| = |V^+|$ (Theorem~\ref{thm:neighbor-fn-bijective}), so $n = 2|V^+|$.
  The restriction of $H$ to the edges between $V^-$ and $V^+$ is a temporal
  bi-clique (Corollary~\ref{cor:biclique-of-not-two-hop}). A one-sided spanner
  of this bi-clique from $V^+$ to $V^-$, of size at most
  $\mathrm{spannerBound}(|V^+|)$
  (Theorem~\ref{thm:exists-spanner-from-to-ncard-le}), together with the two
  perfect matchings $M^-$ and $M^+$ (Lemma~\ref{lem:vminus-vplus-matchings}),
  is a spanner of $H$ (Theorem~\ref{thm:spanner-of-not-two-hop-dismountable}).
\item
  The one-sided bound comes from the halving recurrence of Angrick et al.:
  $\mathrm{spannerBound}(n) \le 2n\lceil\log_2 n\rceil + n$
  (Lemma~\ref{lem:spanner-bound-closed-form}).
\end{enumerate}
